\documentclass[tikz, border=10pt, class=scrbook, 11pt]{standalone}

\usepackage{fontspec}
\defaultfontfeatures{Ligatures=TeX}

\usepackage{amsmath}
\usepackage{amssymb}
\usepackage{mathtools}
\usepackage[
  math-style=ISO,
  bold-style=ISO,
  sans-style=italic,
  nabla=upright,
  partial=upright,
]{unicode-math}
\setmathfont{Latin Modern Math}
\usepackage{optikz}
\usetikzlibrary{shapes.arrows}

\begin{document}
\begin{tikzpicture}[>=stealth]

% ----------------------------------------------------------------------
% Farben: gelb fuer die Praefix-Kaestchen unter der Achse, blau fuer
% die Kaestchen "electronics"/"photonics" oberhalb der Achse.
% Rot (Standardfarbe "red") wird fuer die THz-Luecke und den
% Gammastrahlenbereich verwendet.
% ----------------------------------------------------------------------
\definecolor{praefixgelb}{HTML}{FFF799}
\definecolor{fachblau}{HTML}{639A00}

% ----------------------------------------------------------------------
% Frequenzachse: ein dicker blauer Pfeil, der die Frequenz von 10^0 Hz
% bis 10^24 Hz repraesentiert. Die x-Koordinate ist so skaliert, dass
% eine Einheit einem Sprung von drei Dekaden entspricht -- 10^0, 10^3,
% 10^6, ... liegen also bei x = 0, 1, 2, ..., 8. Das macht die
% Platzierung der folgenden Markierungen sehr einfach.
% ----------------------------------------------------------------------
\draw[very thick, ->, fachblau] (-0.5,0) -- (8.8,0);

% ----------------------------------------------------------------------
% Zehnerpotenzen: eine Schleife ueber alle 9 Markierungen auf der
% Achse. Pro Markierung wird ein kleiner senkrechter Strich auf die
% Achse gezeichnet und die Zehnerpotenz "$10^{e}$" darueber notiert.
% ----------------------------------------------------------------------
\foreach \x/\e in {0/0, 1/3, 2/6, 3/9, 4/12, 5/15, 6/18, 7/21, 8/24} {
    \draw[thick] (\x,-0.1) -- (\x,0.1);
    \node[above] at (\x,0.15) {$10^{\e}$};
}

% ----------------------------------------------------------------------
% SI-Praefixe: dieselben Markierungen bekommen (ausser 10^0, dafuer
% gibt es keinen Praefixnamen) ein gelbes Kaestchen mit dem passenden
% SI-Praefix darunter. Das Praefix "tera" wird dabei rot eingefaerbt,
% weil bei 10^12 Hz genau die THz-Luecke liegt, die weiter unten noch
% gesondert hervorgehoben wird.
% ----------------------------------------------------------------------
\foreach \x/\e/\p in {1/3/kilo, 2/6/mega, 3/9/giga, 4/12/tera,
                      5/15/peta, 6/18/exa, 7/21/zetta, 8/24/yotta} {
    \node[fill=praefixgelb, rounded corners=1pt, inner sep=2pt] at (\x,-0.6) {%
        \ifnum\e=12 \color{red}\fi \p%
    };
}

% ----------------------------------------------------------------------
% "microwaves": Sammelbegriff oberhalb der Achse ueber dem Bereich der
% klassischen Funk-/Mikrowellenbaender, darunter (kleiner gesetzt) die
% Liste der einzelnen Baender von Medium- bis Extremely-High-Frequency.
% ----------------------------------------------------------------------
\node[font=\bfseries] at (2,1) {Mikrowellen};
% \node[font=\small] at (2.2,1.15) {MF, HF, VHF, UHF, SHF, EHF};

% ----------------------------------------------------------------------
% "visible": Beschriftung plus ein kleiner Farbverlaufsbalken, der das
% sichtbare Spektrum stark vereinfacht andeutet (violett - gruen - rot
% statt eines echten Regenbogenverlaufs, um die Zeichnung einfach zu
% halten).
% ----------------------------------------------------------------------
\node[font=\bfseries] at (4.7,1) {Sichtbar};
\shade[left color=red, right color=violet, middle color=green]
    (4.6,-0.2) rectangle (4.9,0.2);

% ----------------------------------------------------------------------
% "x-ray" und "gamma-ray": Beschriftungen fuer die hochenergetischen
% Bereiche jenseits des sichtbaren Lichts, ganz analog zu "microwaves"
% und "visible" einfach oberhalb der Achse platziert.
% ----------------------------------------------------------------------
\node[font=\bfseries] at (6,1.5) {Röntgen};
\node[font=\bfseries] at (7.6,1) {$\gamma$-Strahlung};

% ----------------------------------------------------------------------
% "electronics"/"photonics": zwei blaue Kaestchen ganz oben im Bild.
% "electronics" steht fuer den Bereich, der klassisch elektronisch bis
% kurz vor die THz-Luecke erzeugt/detektiert werden kann, "photonics"
% fuer den Bereich ab der THz-Luecke aufwaerts (Infrarot, sichtbares
% Licht, Roentgenstrahlung), der stattdessen optisch adressiert wird.
% Die Luecke zwischen den beiden Kaestchen liegt bewusst genau ueber
% der 10^12-Hz-Markierung.
% ----------------------------------------------------------------------
\node[fill=fachblau, rounded corners=2pt, text=white, font=\bfseries,
      minimum width=2.6cm, minimum height=0.6cm] at (2.1,2.6) {Elektronische \\ erzeugung};
\node[fill=fachblau, rounded corners=2pt, text=white, font=\bfseries,
      minimum width=2.2cm, minimum height=0.6cm] at (6,2.6) {Schwarzkörper \\ Strahlung};

% ----------------------------------------------------------------------
% THz-Luecke: das eigentliche Herzstueck der Grafik. Ein roter Kreis
% (Ellipse) markiert die 10^12-Hz-Stelle direkt auf der Achse, ein
% roter Pfeil -- dieselbe "single arrow"-Form aus der
% shapes.arrows-Bibliothek, die auch schon in thz_erz.tex fuer die
% Puls-Pfeile verwendet wird, hier um -90 Grad gedreht, damit er nach
% unten zeigt -- verbindet die Beschriftung "THz" mit dieser Stelle.
% Er sitzt bewusst genau in der Luecke zwischen den Kaestchen
% "electronics" und "photonics".
% ----------------------------------------------------------------------
\draw[very thick, red] (4,0) ellipse (0.35 and 0.3);
\node[single arrow, draw=red, fill=red, ultra thick,
      minimum width=0.5cm, minimum height=0.5cm, rotate=-90] at (4,1.65) {};
\node[red, font=\bfseries] at (4,2.1) {THz};

% ----------------------------------------------------------------------
% Achsenbeschriftung unterhalb der gesamten Grafik.
% ----------------------------------------------------------------------
\node[font=\bfseries] at (4,-1.4) {Frequency (Hz)};

\end{tikzpicture}
\end{document}
